\RequirePackage{fix-cm}
\documentclass[10pt]{article}
\usepackage[letterpaper,margin=.38in]{geometry}
\usepackage{amsmath,amssymb,multicol}
\pdfinfoomitdate=1\pdftrailerid{}
\pagestyle{empty}\setlength{\parindent}{0pt}\setlength{\parskip}{2pt}
\setlength{\columnsep}{13pt}\setlength{\jot}{0pt}
\setlength{\abovedisplayskip}{3pt}\setlength{\belowdisplayskip}{3pt}
\let\bigsum\sum\renewcommand{\sum}{\bigsum\nolimits}
\newcommand{\fl}[1]{\left\lfloor#1\right\rfloor}
\newcommand{\rg}[2]{\{#1,\ldots,#2\}}
\begin{document}
\begin{center}
{\large\bfseries A twenty-two-variable candidate equation}\\[-1pt]
{\small Positive integer unknowns; degree 10; a 114-operation evaluator.}
\end{center}\vspace{-8pt}
\begin{multicols}{2}
\fontsize{8.6}{9.1}\selectfont
The only unknowns are the twenty-two distinct positive integers
\[
\begin{gathered}J,F,\alpha,f,g,T,X,Y,\xi,\eta,\zeta,\upsilon,\gamma,\psi,\varphi,\epsilon,\\\ell,n,Z,D,V,\delta\in\mathbb Z_{>0}.\end{gathered}
\]
{\raggedright Their ASCII names in the circuit are respectively
\texttt{J, F, alpha, f, g, T, X, Y, xi, eta, zeta, upsilon, gamma, psi,
varphi, epsilon, ell, n, Z, D, V, delta}. Each name denotes one unknown.\par}
The four fixed integers are \texttt{beta}, \texttt{K}, \texttt{mu}, and
\texttt{omega}, defined below. The single equation is
\[
\boxed{\texttt{hDJ}=0.}
\]
Here \texttt{hDJ} is the last gate of the complete circuit on the next page.
Every \texttt{h}-name denotes an integer polynomial abbreviation, not an
additional unknown. All arithmetic in the circuit is over the integers.

\textbf{The four fixed integers.}
Every index, set, sum, product and floor below defines a constant, not
another unknown. Put
\[
\begin{gathered}
\rho=2^{47},\quad B=\rho^{2011},\quad\beta=B-1,\\
G_m(x)=\sum_{j=0}^{m-1}x^j,\quad
\Lambda=\rho^{1005}.
\end{gathered}
\]
\begin{align*}
\mu={}&\beta-G_{951}(\rho)\\
&-\rho^{960}G_{26}(\rho)-\rho^{1004},\\
\omega={}&\beta+8G_{917}(\rho)-4+(2^{45}-4)\Lambda\\
&+\rho^{960}G_9(\rho)+\rho^{987}G_8(\rho)
  +\rho^{1004},\\
K={}&\rho^{18}+\rho^{1004}+\rho^{2020}\\
&+2\Lambda\Phi(1/\rho),
\end{align*}
where
\[
\begin{gathered}
\chi_i=1+2^{10}(1+w_i+2^8n_i\\
\qquad+2^{17}p_i+2^{25}s_i)\quad(0\le i\le941),\\
\pi(0)=0,\quad\pi(i)=1+\sum_{j=1}^{941}
\fl{\frac{2^{44}-1+\chi_j-\chi_i}{2^{44}}}\quad(1\le i\le941).
\end{gathered}
\]
\[
\begin{split}
\Phi(x)={}&\sum_{i=0}^{941}\chi_i x^{\pi(i)}\\
&+2^{10}x^{969}(x^9+2^{17})G_8(2x)\\
&+2^{18}x^{942}(x^{18}+2^{17})G_9(2x).
\end{split}
\]
The rational intermediate in K cancels exactly; all four results are positive integers.

\textbf{The 942 indexed fixed tiles.}
Use the actions $u_j,v_j,d_j$ decoded below. For $d=0,1$, let
\[
\begin{gathered}
h(q,b)=2q+b+2,\quad
\mathcal S_d=\{u_j:d_j=d,\ 0\le j<430\},\\
r_{d,q}=205-199d+|\mathcal S_d\cap\rg0{q-1}|.
\end{gathered}
\]
An interval ending at $-1$ is empty. The table defines $t_i$ for
$0\le i\le941$; all index ranges are fixed:
\[
\setlength{\arraycolsep}{1.4pt}\begin{array}{c|l}
\text{index }i&\text{tile }t_i=(w_i,n_i,p_i,s_i)\\\hline
0&(2,435,3,1)\\
j+1\ (j=2q+b<430)&((1-d_j)r_{0,u_j},h(q,b),d_jr_{1,u_j},v_j)\\
433-2a+b\ (a,b=0,1)&(0,432a+b,0,432a+b)\\
\substack{423+2r_{d,q}+b\\q\in\mathcal S_d,\ d,b=0,1}&(dr_{d,q},b,(1-d)r_{d,q},h(q,b))\\
923+2a+b\ (a,b=0,1)&(0,432a+b,0,438)\\
927&(2,434,2,0)\\
928&(3,436,4,4)\\
929&(4,436,4,0)\\
496+t\ (t=434,435,436)&(1,438,1,t)\\
933+3a+b\ (a,b=0,1,2)&(a^2,437,b,437)
\end{array}
\]
\textbf{A closed decoder for the 430 actions.}
Put $D_b(A,k)=\fl{A/b^k}\bmod b$, where
$a\bmod b=a-b\fl{a/b}$. For $0\le q<215$, $b=0,1$, $j=2q+b$:
\[
\begin{gathered}
\mathcal N=\mathcal A\bmod2^{636},\quad
\mathcal T=\fl{\mathcal A/2^{636}},\\
\nu_j=D_2(\mathcal N,j),\quad \kappa_j=(\sigma_j+2)^{1-\nu_j}\quad(0\le j<430),\\
\sigma_j=\sum_{k=0}^{j-1}\nu_k,\quad
\Pi_j=\prod_{k=0}^{j-1}\kappa_k\quad(0\le j\le430),\\
m_q=D_2(\mathcal N,430+q),\quad \tau_q=\sum_{k=0}^{q-1}m_k,\\
F_q=14+m_q(D_{16}(\fl{\mathcal T/\Pi_{430}},\tau_q)-14),\\
u_j=(\sigma_{j+1}-(\fl{\mathcal T/\Pi_j}\bmod\kappa_j))\bmod216,\\
v_j=D_2(F_q,2b),\quad d_j=D_2(F_q,2b+1).
\end{gathered}
\]
Empty sums/products are $0/1$. Here $(v_0,d_0,u_0)=(1,1,1)$; the constructor tile
uses this actual first target, including its state renaming.

\textbf{The integer $\mathcal A$ in base 10.}
Its complete decimal expansion is the concatenation of these lines:\par
\begingroup\fontsize{7.7}{8.2}\selectfont\ttfamily
12889260410359561746736540743515845134244260640012702033212621\par
95398074180784058571518488282119861023093355389406939414042892\par
97455773992193784795705001618777025966369137954859764532832722\par
75411843410809038713502633204617333246733616940964578272360901\par
56186216508504953649192752154519268024509952082950238865468307\par
22999632482769253980274091217460707829023970817244790340704449\par
09085581271192071656170732004937337082171141082486384269072884\par
72778087419663417942892520175630649792329265834490253035577533\par
96449689808375788457488971017043689875433162972390177936123032\par
89025657757489526342497438772045714099661532481123050235898973\par
27443609903800541199351405985754278307227777584454508422717998\par
334113564409335346074892680328437\par
\endgroup
\normalfont\fontsize{7.7}{8.2}\selectfont
The fixed integer $\mathcal A$ above is written in base 10. This is a research
candidate. Exact polynomial calculations certify degree 10 and the
substitution identities. Historical paired Lean runs compiled the witness
transport and an internal intuitionistic PRA equivalence to a translated
consistency sentence. No
minimality claim is made.
\end{multicols}
\clearpage
\begin{center}{\large\bfseries Complete 114-operation circuit}\end{center}
Each line defines one gate. The only inputs are the 22 positive integers and
the four fixed integers on the first page. The final gate \texttt{hDJ} is a
sum of 13 squares; setting it to zero is the displayed equation. Gate names
are local abbreviations and are distinct from all input names.
\begin{multicols}{2}
\raggedright\ttfamily\fontsize{8}{9.2}\selectfont
% BEGIN CIRCUIT
1. hA = J * beta\\
2. hB = 1 + hA\\
3. hC = hB * hB\\
4. hD = hB * hC\\
5. hE = Y * hC\\
6. hF = hB - F\\
7. hG = hF - alpha\\
8. hH = X * Y\\
9. hI = eta + zeta\\
10. hJ = Y * hI\\
11. hK = eta + hJ\\
12. hL = hK * hK\\
13. hM = hH + Y\\
14. hN = 4 * hM\\
15. hO = 3 + hN\\
16. hP = hM * hM\\
17. hQ = hO + hP\\
18. hR = hM * hK\\
19. hS = X + hR\\
20. hT = hH * hJ\\
21. hU = hD * hA\\
22. hV = K * hB\\
23. hW = X + hV\\
24. hX = hG * hW\\
25. hY = F * hB\\
26. hZ = hX - hY\\
27. hAA = ell + hI\\
28. hAB = hAA + xi\\
29. hAC = hI * hZ\\
30. hAD = hAB - hAC\\
31. hAE = psi * hAD\\
32. hAF = Y * hA\\
33. hAG = hZ * hD\\
34. hAH = hAE + hAF\\
35. hAI = hAH - hAG\\
36. hAJ = Z * Z\\
37. hAK = psi * psi\\
38. hAL = ell * hAA\\
39. hAM = 1 + hAL\\
40. hAN = hAK * hAM\\
41. hAO = hAJ - hAN\\
42. hAP = hF * hA\\
43. hAQ = alpha + hAP\\
44. hAR = hC - 1\\
45. hAS = hAQ * hAR\\
46. hAT = omega * hB\\
47. hAU = mu + hAT\\
48. hAV = J * hAU\\
49. hAW = hAS + hAV\\
50. hAX = hI - epsilon\\
51. hAY = hAX - 1\\
52. hAZ = varphi * gamma\\
53. hBA = hH * hB\\
54. hBB = hAZ + hBA\\
55. hBC = hBB - hAY\\
56. hBD = varphi * hS\\
57. hBE = hO * hBC\\
58. hBF = hBD + hBE\\
59. hBG = D * D\\
60. hBH = varphi * varphi\\
61. hBI = delta * n\\
62. hBJ = 1 + hBI\\
63. hBK = hBH * hBJ\\
64. hBL = hBG - hBK\\
65. hBM = f * f\\
66. hBN = hBM - 1\\
67. hBO = delta * hBN\\
68. hBP = n * g\\
69. hBQ = hBP * hBP\\
70. hBR = hBQ - hBO\\
71. hBS = T * hK\\
72. hBT = hBS - epsilon\\
73. hBU = hBM * hBT\\
74. hBV = hBU - hK\\
75. hBW = upsilon * upsilon\\
76. hBX = V * V\\
77. hBY = hBW - hBX\\
78. hBZ = hBO * hBY\\
79. hCA = hBW - hBZ\\
80. hCB = hCA - 1\\
81. hCC = hAO * hAO\\
82. hCD = hBL * hBL\\
83. hCE = hCB * hCB\\
84. hCF = hBR * hBR\\
85. hCG = hCC + hCD\\
86. hCH = hCG + hCE\\
87. hCI = hCH + hCF\\
88. hCJ = psi - hU\\
89. hCK = hCJ * hCJ\\
90. hCL = hCI + hCK\\
91. hCM = varphi - hE\\
92. hCN = hCM * hCM\\
93. hCO = hCL + hCN\\
94. hCP = epsilon - hAW\\
95. hCQ = hCP * hCP\\
96. hCR = hCO + hCQ\\
97. hCS = ell - hT\\
98. hCT = hCS * hCS\\
99. hCU = hCR + hCT\\
100. hCV = n - hL\\
101. hCW = hCV * hCV\\
102. hCX = hCU + hCW\\
103. hCY = Z - hAI\\
104. hCZ = hCY * hCY\\
105. hDA = D - hBF\\
106. hDB = hDA * hDA\\
107. hDC = V - hBV\\
108. hDD = hDC * hDC\\
109. hDE = delta - hQ\\
110. hDF = hDE * hDE\\
111. hDG = hCX + hCZ\\
112. hDH = hDG + hDB\\
113. hDI = hDH + hDD\\
114. hDJ = hDI + hDF\\

% END CIRCUIT
\end{multicols}
\end{document}
